% ===== 原文 pp.43–44（含 p.44 末句跨页补完至 p.45 首行） =====
% 译者 tr-06：§5.2 后半：斜率与稳定性、多边形/超多边形空间、图 5.10、定理 5.11、注 5.12
% 编号区间：图 5.10、定理 5.11、注 5.12
% 依据：split/p043.txt、split/p044.txt、split/p045.txt 首行（仅用于定位）＋ pages/p043.png、p044.png（200 dpi，数学逐符号核对）

% [原文 p.43]
与处理抛物 Higgs 丛时类似，人们可以定义箭图的斜率与稳定性条件。对稳定性参数 $\alpha := (\alpha_v)_{v\in V}\in\Z^V$，$W$ 的 $\alpha$-斜率为
\[
\mu_{\alpha}(W):=\frac{\sum_{v\in V}\alpha_v\dim W_v}{\sum_{v\in V}\dim W_v}\in\Q.
\]
于是称表示 $W$ 为：$\alpha$-半稳定的，若对所有子表示 $0\neq W'\subset W$ 都有 $\mu_{\alpha}(W')\leqslant\mu_{\alpha}(W)$；$\alpha$-稳定的，若对所有子表示 $0\neq W'\subsetneq W$ 都有 $\mu_{\alpha}(W')<\mu_{\alpha}(W)$。相应地有自然的 Jordan--Hölder 滤链概念（如第 2.9 节所考虑的那些），并且若两个 $\alpha$-半稳定表示关于其 Jordan--Hölder 滤链的相伴分次对象同构，则称它们 S-等价。更多细节，读者可参阅例如 [66]，本节大部分记号遵循该文。

% [Tutor Note: p.43 原文 "（as those considered in Section 2.9）" 有误：本讲义只有 2.1 至 2.5 五个小节，不存在第 2.9 节；Jordan--Hölder 滤链的实际出处是 §2.1（p.4，式 (2.7) 之前那一段：$F_0=0\subset F_1\subset\dots\subset F_r=E$ 与 $Gr(E,\Phi):=\bigoplus_i(F_i/F_{i-1},\Phi_i)$，紧接着定义 S-等价）。故 "Section 2.9" 应为 "Section 2.1"。译文照原样排 2.9。]

我们对箭图簇的兴趣来自它们与向量丛、Higgs 丛的关系，这种关系特别地可以通过考虑多边形空间与超多边形空间而得到（例如见 [50] 及其中的参考文献）。

多边形空间 $\cM_{\mathrm{poly}}(n,\ell)$ 由欧氏空间 $\E^3$ 中长度由 $\ell=(\ell_1,\dots,\ell_n)$ 给出的 $n$ 边形组成，模去保向等距作用；它同胚于 $\PP^1$ 上 $n$ 个有序点的 $\ell$-半稳定模空间，再模去射影直线的自同构（见 [66] 及其中的参考文献）。此外，它也同构于某个星形箭图 $Q^{\ell}_n$ 的表示模空间。

% [Tutor Note: p.43 原文 "is homeomorphic the moduli space" 漏了介词，应为 "is homeomorphic to the moduli space"；译文按正常语义处理。]

为了定义超多边形空间的模空间，我们考虑箭图 $Q=(V,A,h,t)$，并定义加倍箭图（doubled quiver）为 $\overline{Q} := (V,\overline{A},h,t)$，其中 $\overline{A}=A\cup A^*$，而
\[
A^* := \{a^* : h(a)\to t(a)\}_{a\in A}.
\]

通过星形箭图的加倍箭图，[42] 的作者对任意的 $r$ 构造了超多边形空间；它们是多边形空间的超凯勒类比（关于以较不一般的形式出现的早期工作，见该文中的参考文献）。

\begin{figure}[htbp]
\centering
\includegraphics[width=0.70\linewidth]{fig2.png}
\figtitle{5.10}{左：其表示模空间为多边形空间的箭图；右：用于构造超多边形空间的加倍箭图。}
\end{figure}

% [原文 p.44]
延续 [42] 的工作，通过群 $G := \mathrm{P}(\U(r)\times(S^1)^n)$ 对图 5.10 中星形箭图的自然作用，$G$ 在 $T^*\Rep(Q)$ 上的作用有两个自然的矩映射，其标量倍数可写成如下形式：
\[
\text{实矩映射}\quad \mu_r(x,y)=((xx^*-y^*y)_0,\ |x_1|^2-|y_1|^2,\dots,|x_n|^2-|y_n|^2)
\]
\[
\text{复矩映射}\quad \mu_c(x,y)=((xy)_0,\ y_1x_1,\dots,y_nx_n)
\]
其中 $(\cdot)_0$ 表示矩阵的无迹部分。多边形空间 $\cP^r_n(\alpha)$ 是辛商
\[
\cP^r_n(\alpha)=\Rep(Q)/\!/_{\alpha}G.
\]
类似地，超多边形空间 $\cX^r_n(\alpha)$ 可定义为超凯勒商
\[
\cX^r_n(\alpha)=T^*\Rep(Q)/\!/\!/_{(\alpha,0)}G.
\]

% [Tutor Note: p.44 双斜杠 $/\!/$ 表示辛商（symplectic quotient），三斜杠 $/\!/\!/$ 表示超凯勒商（hyperkähler quotient）：后者用实、复、实三个矩映射作商，原文把剩余参数写成下标 $(\alpha,0)$，即实矩映射取 $\alpha$、复矩映射取 $0$。]
% [Tutor Note: p.44 矩映射的取值可核对：$\Rep(Q)=\C^{n\times r}$，故 $T^*\Rep(Q)$ 的坐标为 $(x,y)$，$x$ 是 $n\times r$ 矩阵、$y$ 是 $r\times n$ 矩阵；于是 $xx^*-y^*y$ 为 $r\times r$ 矩阵（取迹零部分 $(\cdot)_0$，取值于 $\fsu(r)$），而 $|x_i|^2-|y_i|^2$ 与 $y_ix_i$ 各给出 $n$ 个实/复标量，恰好对偶于 $G=\mathrm{P}(\U(r)\times(S^1)^n)$ 的 Lie 代数分解 $\fsu(r)\oplus\R^n$，这正是原文 "two natural moment maps ... whose scalar multiples" 中「两个」矩映射的来由。]
% [Tutor Note: p.44 原文 "can be defined be the hyperkähler quotient" 中第二个 "be" 应为 "by"（原文笔误）；同页 "on the star-shaped quivers in Figure 5.10" 用复数 "quivers"，而图 5.10 给的是左（箭图）右（加倍箭图）两个，复数可以成立，译文照原样。]

Nakajima 证明了多边形空间与超多边形空间分别是凯勒与超凯勒箭图簇的例子，它们随后可用来理解抛物 Higgs 丛。尽管箭图簇与抛物 Higgs 丛之间的关系已被长期研究，但仍有许多有趣的问题尚未解决。

基于复射影空间中半稳定 $n$ 边形的模空间的超凯勒版本与某些抛物 Higgs 丛之间的自然类比，Rayan 与 Fisher 在 [42] 中把空间 $\cX^r_n(\alpha)$ 等同于退化 Hitchin 系统：对一个超多边形 $(x,y)\in\cX^r_n(\alpha)$，以及 $\Sigma$ 上固定的除子 $D=\sum_{i=1}^{n}p_i$（其中任两个 $p_i$ 都不相同），$\PP^1$ 上以向量丛 $E=\C^r\times\PP^1$ 为底丛的 Higgs 场由
\[
\phi(z)=\sum_{i=1}^{n}\frac{\phi_i\dd z}{z-p_i},
\]
给出，其中 $z$ 是 $\PP^1$ 上的仿射坐标，且 $\phi_i=x_iy_i$。

更确切地说，对抛物 $\SL(n,\C)$-Higgs 丛，有如下一般结果（关于抛物 Higgs 丛如何对应于超多边形，[42, 第 4 节] 给出了清楚的解释）：

\begin{thm}{5.11（[42]）}
空间 $\cX^r_n(\alpha)$ 可等同于 $\PP^1$ 上秩为 $r$ 的抛物 Higgs 丛的模空间，使得每个标记点处的留数落在 $\fsl_r$ 的极小幂零轨道的闭包中，且其底丛平凡。
\end{thm}

此外，在上述等同之下，如 [42] 所示，存在一个自然的 Hitchin 映射，其底空间是一个仿射空间，其维数为超多边形空间总维数的一半。

\begin{rem}{5.12}
上述一般超多边形空间的数学构造应当会在弦论中引出有趣的应用。特别地，由 [4]，这些箭图可由 F-理论中的一个相交膜组态生成。中心节点对应于来自缠绕在规范场论曲线上的 7-膜的贡献，而卫星节点对应于系统的味膜。
\end{rem}

% 跨页句说明：p.44 末句 "The central node corresponds to the contribution from the 7-brane" 与
% p.45 首行 "wrapped over the gauge theory curve, and the satellite nodes to the flavour branes of the system."
% 是同一句，按 TIMELINE #7 由 tr-06 在本块补完（"wrapped over the gauge theory curve" 是 "the 7-brane" 的后置修饰语）。
% tr-07 自 p.45 的 §5.3 Langlands 对偶起写，不得重复此句。
